% answers_short.tex —— 2010 年·普通高等学校招生考试·四川卷·文科 · 简短答案速查
% 由 tools/make_short_answers.py 自 answers.tex 的 \gksolhead/\gkanswer 提取，
% 并与 2010/2010_数学真题_四川文_答案要点.md 逐题比对后生成，勿手改（改源 MD 后重跑）。

\begin{gknotice}
\textbf{使用说明}：本节只给最终答案便于快速核对，完整推导见解析版。答案为编写者独立推导，\textbf{非官方标准答案}。
\end{gknotice}

\gkpart{选择题}
\begin{center}\begin{tabular}{ccccccccccccc}
\toprule
题号 & 1 & 2 & 3 & 4 & 5 & 6 & 7 & 8 & 9 & 10 & 11 & 12 \\
\midrule
答案 & D & C & C & D & A & C & C & B & A & D & D & A \\
\bottomrule
\end{tabular}\end{center}

\gkpart{填空题}
\begin{itemize}[leftmargin=2.6em]
  \item[\textbf{13.}] $24$。
  \item[\textbf{14.}] $2\sqrt3$。
  \item[\textbf{15.}] $\dfrac{\sqrt3}4$。
  \item[\textbf{16.}] ①②。
\end{itemize}

\gkpart{解答题}
\begin{itemize}[leftmargin=2.6em]
  \item[\textbf{17.}] （1）$\dfrac{125}{216}$；（2）$\dfrac{25}{27}$。
  \item[\textbf{18.}] （1）证明；（2）$\arccos\dfrac13$。
  \item[\textbf{19.}] （1）证明；（2）$\dfrac{3\sqrt{10}}{10}$。
  \item[\textbf{20.}] （1）$a_n=4-n$；（2）$q=1$ 时 $S_n=\dfrac{n(n+1)}2$，$q\neq1$ 时 $S_n=\dfrac{1-(n+1)q^n+nq^{n+1}}{(1-q)^2}$。
  \item[\textbf{21.}] （1）$x^2-\dfrac{y^2}3=1$；（2）过点 $F$。
  \item[\textbf{22.}] （1）$g(x)=\log_a\dfrac{x-1}{x+1}$；（2）$a>1$ 时 $0<t<5$，$0<a<1$ 时 $t>32$；（3）$f(1)+\cdots+f(n)<n+4$。
\end{itemize}
